\documentclass{situationist-preprint}

\title{Partiality Without One Thought Too Many}
\msid{ETH-19-a}
\msfamily{ETH-19}
\msdiscipline{Normative and Applied Ethics}
\mstype{Principal manuscript}
\mskeywords{partiality; impartiality; one thought too many; special obligations; relationships; agent-relative reasons; personal prerogative}

\begin{document}
\maketitle

\begin{abstract}
A man who can save his wife or a stranger, but not both, may save his wife. Impartialist theories accept the verdict but hold that it needs an impartial justification: it must be shown that a practice of partiality, or a disposition to it, is what an impartial standpoint would endorse. Williams objected that a man who thinks this has \enquote{one thought too many}. We argue that Williams was right, and that agents may give priority to those close to them without justifying this priority by impartial reasons. The argument has two parts. The first is a modal argument: every impartial justification makes the permission depend on contingent facts, such as the welfare effects of partial practices, on which the permission does not in fact depend. The second is an argument from the value of relationships: to value a relationship non-instrumentally is to treat it as a source of reasons, so that a relationship whose reasons required impartial certification would not be valued as relationships are. Partiality is limited by the duties of beneficence and rescue, but limited is not the same as justified, and above those duties the personal prerogative may be spent on those close to one.
\end{abstract}

\section{The problem}

\citet{Godwin1793} considered a case in which a fire breaks out in the palace of Archbishop F\'enelon, and one can save either the archbishop, whose work will benefit thousands, or his chambermaid, who may be one's mother. He held that one ought to save the archbishop, and asked what there could be in the pronoun \emph{my} that should overturn the decisions of impartial truth. The modern form of the case is due to \citet{Fried1970}, who imagined a man who can rescue his wife or a stranger, and to \citet{Williams1981}, who argued that if the man's justification for saving his wife were that it is permissible, in situations of this kind, to save one's wife, he would have \enquote{one thought too many}: it might have been hoped that his motivating thought would be that it was his wife.

The literature since has divided into three positions.
\begin{enumerate}[label=(\roman*)]
\item \emph{Strict impartialism}: partiality is not justified where it departs from what impartial reasons support. Godwin held this view, and it is the view of naive act consequentialism.
\item \emph{Impartial justification of partiality}: partiality is justified, but its justification is impartial. Rule consequentialists hold that a code permitting partiality has the best consequences \citep{Hooker2000}; sophisticated consequentialists hold that dispositions to partiality are optimific, even if acts flowing from them are not always \citep{Railton1984}; some Kantians hold that partial acts are permissible because a maxim permitting them could be universally willed \citep{Baron1991}.
\item \emph{Non-derivative partiality}: relationships give reasons of their own, which do not need impartial justification \citep{Williams1981,Scheffler1997,Kolodny2003,Keller2013}.
\end{enumerate}

We defend (iii). The defenders of (ii) agree that partiality is permissible and disagree only about why. The disagreement is not idle: it determines whether the permission would survive a change in the facts on which the impartial justification relies, and it determines what the agent who acts well must think.

\section{Framework}

\begin{definition}[Partial reason]\label{d:partial}
A \emph{partial reason} is a reason for an agent to give priority to the interests of a person $Y$ that is grounded in a relationship between the agent and $Y$, such as kinship, friendship, love or membership of a common association.
\end{definition}

\begin{definition}[Impartial justification]\label{d:impartial}
An \emph{impartial justification} of a permission to act on partial reasons is a derivation of the permission from a principle that gives no fundamental weight to the identity of the agent or to her relationships, together with facts about the effects or the universal adoptability of acting partially.
\end{definition}

\begin{principle}[Modal robustness]\label{p:robust}
If a permission is grounded in some facts, then it varies with those facts: in circumstances in which the grounding facts fail, and nothing else grounds the permission, the permission fails.
\end{principle}

\begin{principle}[Constitution]\label{p:constit}
To value a relationship non-instrumentally is, in part, to treat it as a source of reasons for action that are not derived from the value of anything outside it.
\end{principle}

Principle~\ref{p:robust} is a general constraint on grounding and is accepted by the defenders of (ii), who take the permission to be grounded in impartial facts. Principle~\ref{p:constit} is the account of valuing a relationship given by \citet{Scheffler1997} and developed by \citet{Kolodny2003}; we defend it in Section~4.

\section{The argument}

\begin{lemma}\label{l:contingent}
Every impartial justification of a permission to act on a partial reason grounds the permission in contingent facts about the effects or the general adoptability of partial conduct.
\end{lemma}

\begin{proof}
By Definition~\ref{d:impartial}, the principle from which the permission is derived gives no fundamental weight to the agent's relationships. The permission must therefore be derived from facts about what partial conduct brings about or about what would follow if everyone engaged in it: that partial dispositions sustain close relationships and that close relationships promote welfare \citep{Railton1984}, that a code permitting partiality has lower costs of internalization than one forbidding it \citep{Hooker2000}, or that a world of universal partiality is one an agent could will. These are facts about human psychology and social organization, which could have been otherwise.
\end{proof}

\begin{lemma}\label{l:invariant}
The permission to give priority to those close to one does not vary with the contingent facts of Lemma~\ref{l:contingent}.
\end{lemma}

\begin{proof}
Consider a world in which human psychology is such that partial dispositions produce somewhat less welfare than impartial ones, because, say, impartial carers are slightly more efficient, and in which a code forbidding partiality would be easy to internalize. In that world, Fried's man faces the same choice between his wife and a stranger. The impartial justifications of (ii) imply that, in that world, he may not save his wife, or may save her only if the welfare difference happens to be small. We judge that he may save her in that world as in ours. The judgment that he may save his wife is not hostage to the efficiency of carers. It is a judgment about what his relationship to her permits him to do, and the relationship is the same in both worlds.
\end{proof}

\begin{theorem}\label{t:main}
Agents may give priority to those close to them, and the permission is not grounded in an impartial justification.
\end{theorem}

\begin{proof}
That the permission exists is common ground between positions (ii) and (iii), and is supported by the considered verdicts in Fried's and Godwin's cases. Suppose it were grounded in an impartial justification. By Lemma~\ref{l:contingent} it would then be grounded in contingent facts about the effects or adoptability of partial conduct, and by Principle~\ref{p:robust} it would fail in circumstances in which those facts fail. By Lemma~\ref{l:invariant} it does not fail in such circumstances. So it is not grounded in an impartial justification. What grounds it is the relationship itself, which supplies partial reasons in the sense of Definition~\ref{d:partial}.
\end{proof}

\begin{corollary}\label{c:thought}
The agent who saves his wife because she is his wife, without the further thought that partiality is impartially permitted, acts on a sufficient reason.
\end{corollary}

\begin{proof}
By Theorem~\ref{t:main}, the permission is grounded in the relationship. The thought \enquote{she is my wife} therefore picks out the fact that grounds the permission, and nothing further is needed for the agent's reason to be sufficient. The additional thought that the impartial standpoint would endorse his act adds no ground; it is the thought too many.
\end{proof}

\section{The value of relationships}

Theorem~\ref{t:main} rests on Lemma~\ref{l:invariant}, a judgment about a counterfactual case. Principle~\ref{p:constit} gives a second, independent route to the same conclusion, and explains the judgment.

Suppose that a person values her friendship with $Y$ but that she takes the reasons it gives her to be derived from impartial considerations, so that she would treat them as void if those considerations ceased to support them. Then what she values is not the friendship but the impartial good it happens to serve, and her attitude to the friend is that of an agent who would transfer her concern to any relationship that served that good as well. This is not how friends value their friends. As \citet{Stocker1976} argued, a person who visits a friend in hospital because it is his duty, or because it maximizes the good, and not because of the friend, exhibits a split between his reasons and his motives that is a defect in the friendship. If valuing a relationship is constituted by treating it as a source of non-derivative reasons, as Principle~\ref{p:constit} states, then whoever values relationships as they are valued already accepts partial reasons that need no impartial justification. This fits the result of EXI-16 that love responds to reasons that do not make it transferable to qualitatively similar persons: the reasons of love are reasons about this person and this relationship.

The value of relationships is itself one of the components of the good (ETH-01), and an impartialist may reply that it is this value, impartially weighed, that justifies partiality. The reply confuses two things. That relationships are good is a fact the impartial standpoint can register; but the reasons a relationship gives to the people in it are reasons for them to favor each other, not reasons for everyone to promote relationships in general. An impartial derivation can deliver only the second kind of reason, and the second does not entail the first: a reason to promote friendship in general may require one to neglect one's friend in order to foster several friendships elsewhere.

\section{The limits of partiality}

That partiality needs no impartial justification does not mean that it is unlimited. The man may save his wife rather than a stranger; he may not kill a stranger to give his wife a kidney. The limits come from the claims of others, and they include the duties of beneficence and rescue. ETH-06 argues that affluent agents must give a substantial fixed share of their income, ETH-18 that there is a duty to rescue up to moderate cost, and ETH-28 that the combined requirement leaves at least four fifths of income to a personal prerogative. Partial reasons do not override these duties, so that the share required by beneficence may not be redirected to one's family. Within the prerogative, however, the agent may spend her resources favoring those close to her, and Theorem~\ref{t:main} implies that she needs no impartial reason for doing so.

The distinction between limiting and justifying is the key to the position. An impartial constraint that bounds what partial reasons may permit is not a justification of what they permit within the bound. Compare the prerogative to pursue one's own projects: that one may not pursue them at the cost of a child drowning in front of one does not show that one's reason to pursue them is impartial.

\section{Objections and replies}

\begin{objection}[Arbitrariness]\label{o:arb}
Without an impartial justification, partiality is arbitrary. The wife and the stranger have equal claims, and preferring one because of a relation to oneself is like preferring one because of the color of her hair.
\end{objection}

\begin{reply}
The color of a person's hair gives no reason to favor her; a relationship does. The objection assumes that only impartial considerations are reasons, which is what is in question. That a relationship is a reason for those in it is not arbitrary in any sense that should trouble us: it is a reason that anyone in such a relationship has, so that the principle granting it is universal in form, though agent-relative in content.
\end{reply}

\begin{objection}[Bad relationships]\label{o:bad}
If relationships give reasons of their own, then so do relationships of shared racial identity or membership in a criminal gang. A theory of non-derivative partiality must either accept this or appeal to the value of the relationship, which reintroduces an impartial standard.
\end{objection}

\begin{reply}
We accept that only valuable relationships give partial reasons; a relationship constituted by a wrongful purpose gives none. The condition is not an impartial justification in the sense of Definition~\ref{d:impartial}, since it does not derive the permission from effects or adoptability. It is a condition on what counts as a relationship of the relevant kind, in the way that a promise extracted by fraud does not bind without its showing that promissory obligation has an impartial ground.
\end{reply}

\begin{objection}[Sophisticated consequentialism]\label{o:railton}
Railton's sophisticated consequentialist does not have the thought too many. He saves his wife because she is his wife; the consequentialist justification operates at the level of his dispositions, not in his deliberation \citep{Railton1984}.
\end{objection}

\begin{reply}
Railton's view avoids the thought too many in deliberation, and to that extent improves on its predecessors. It does not avoid the dependence of the permission on contingent facts, and the modal argument of Lemma~\ref{l:invariant} applies to it unchanged. It also introduces a gap between what the agent thinks justifies his act and what in fact justifies it, so that the agent who acts well is mistaken about why he acts well. Our view closes the gap: what the man thinks, that she is his wife, is what justifies him.
\end{reply}

\begin{objection}[Godwin's archbishop]\label{o:godwin}
In Godwin's case, the archbishop's work will benefit thousands. Surely one should save him rather than one's mother, which shows that impartial considerations can override partial ones and so must be weighed against them on a common scale.
\end{objection}

\begin{reply}
Partial reasons can be outweighed, and when the stakes for others are great enough they are. This shows that partial and impartial reasons can conflict and that conflicts must be resolved, not that partial reasons derive from impartial ones. Weighing two reasons does not require that one be derived from the other.
\end{reply}

\section{Conclusion}

Partiality is justified, and its justification is not impartial. Every impartial justification makes the permission to favor those close to one depend on contingent facts on which it does not depend, and to value a relationship is to take it as a source of reasons of its own. The man who saves his wife because she is his wife has exactly the thought he needs. Partial reasons are limited by the duties of beneficence and rescue, but within the personal prerogative they may be acted on without further warrant.

\begin{thebibliography}{99}\small
\bibitem[Baron(1991)]{Baron1991} Baron, M. (1991). Impartiality and friendship. \emph{Ethics}, 101(4), 836--857.
\bibitem[Fried(1970)]{Fried1970} Fried, C. (1970). \emph{An Anatomy of Values: Problems of Personal and Social Choice}. Harvard University Press.
\bibitem[Godwin(1793)]{Godwin1793} Godwin, W. (1793). \emph{An Enquiry Concerning Political Justice}. London: G.~G.~J. and J. Robinson. Book II, ch.~2.
\bibitem[Hooker(2000)]{Hooker2000} Hooker, B. (2000). \emph{Ideal Code, Real World: A Rule-Consequentialist Theory of Morality}. Oxford University Press.
\bibitem[Keller(2013)]{Keller2013} Keller, S. (2013). \emph{Partiality}. Princeton University Press.
\bibitem[Kolodny(2003)]{Kolodny2003} Kolodny, N. (2003). Love as valuing a relationship. \emph{The Philosophical Review}, 112(2), 135--189.
\bibitem[Railton(1984)]{Railton1984} Railton, P. (1984). Alienation, consequentialism, and the demands of morality. \emph{Philosophy \& Public Affairs}, 13(2), 134--171.
\bibitem[Scheffler(1997)]{Scheffler1997} Scheffler, S. (1997). Relationships and responsibilities. \emph{Philosophy \& Public Affairs}, 26(3), 189--209.
\bibitem[Stocker(1976)]{Stocker1976} Stocker, M. (1976). The schizophrenia of modern ethical theories. \emph{The Journal of Philosophy}, 73(14), 453--466.
\bibitem[Williams(1981)]{Williams1981} Williams, B. (1981). Persons, character and morality. In \emph{Moral Luck: Philosophical Papers 1973--1980}, Cambridge University Press, 1--19.
\end{thebibliography}

\end{document}
